\section{Three Parallel Theories of Cyclic Torsors}
\label{section:cyclic_extension_theories}
\providecommand{\bW}{\mathbb{W}}

The first reduction of the preceding section explains why cyclic theories
suffice: it reduces the theorem from finite abelian groups to cyclic groups of
order prime to $p$ and cyclic groups of order $p^r$. The second reduction
explains exactly what these theories must accomplish: they must realize a local
cyclic torsor over $k((x))$ by an explicit equation on
$\bG_m\subset\bP^1_k$, with the required control of its ramification. We
therefore recall the three theories in the parallel form in which they will be
used:
\[
\begin{array}{c|c|c|c}
\text{theory}&H&\varphi:H\to H&\ker(\varphi)\\ \hline
\text{Kummer}&\bG_m&z\mapsto z^e&\mu_e\cong\Z/e\Z\\
\text{Artin--Schreier}&\bG_a&z\mapsto z^p-z&\F_p\cong\Z/p\Z\\
\text{Artin--Schreier--Witt}&\bW_r&F-\mathrm{id}
    &W_r(\F_p)\cong\Z/p^r\Z.
\end{array}
\]
Here the identification in the Kummer row is made after choosing a primitive
$e$-th root of unity. In every row, pulling back $\varphi$ along a section $a$
produces the torsor given by the equation $\varphi(X)=a$, and changing $a$ by
an element of $\varphi(H)$ does not change its class. The kernel action,
pullback construction, and resulting cohomology class are detailed in
\Cref{appendix:common_group_scheme_construction}.

Artin--Schreier--Witt theory contains Artin--Schreier theory as the case
$r=1$. We nevertheless treat the Artin--Schreier case separately: it displays,
without the additional Witt-vector bookkeeping, the basic principle later
used in the general case---choose a reduced representative of the local
torsor, control its pole order by the ramification break, and use the same
equation to realize the torsor on $\bG_m$.

Only the conclusions used in Chapter~2 are stated here. The group-scheme
constructions, classification arguments, normal-form results, and proofs are
collected in \Cref{appendix:cyclic_torsor_theories}.

Throughout this section, let $K$ be a complete discrete valuation field with
perfect residue field $\kappa$ of characteristic $p>0$.

\subsection{Kummer theory: order prime to \texorpdfstring{$p$}{p}}

\paragraph{Classes and equations.}
Assume that $K$ contains $\mu_n$ and that $(n,p)=1$. Write
\[
H^1(K,\mu_n):=H^1_{\mathrm{\acute et}}(\Spec K,\mu_n);
\]
equivalently, this is the first continuous Galois cohomology group of $K$
with coefficients in $\mu_n$. Its elements classify $\mu_n$-torsors over
$K$. The Kummer sequence gives a natural isomorphism
\[
H^1(K,\mu_n)\cong K^\times/(K^\times)^n.
\]
Under this isomorphism, the coset of $a\in K^\times$ corresponds to the
torsor
\[
\Spec K[T]/(T^n-a),
\]
on which $\mu_n$ acts by $T\mapsto\zeta T$. In particular, replacing $a$
by $ab^n$, with $b\in K^\times$, does not change the isomorphism class of
the torsor. If the class of $a$ has order $n$, adjoining a solution of
$T^n=a$ gives a cyclic field extension of degree $n$; conversely, every
cyclic extension of degree $n$ arises in this way.

\begin{theorem}[Ramification in Kummer Extensions,
{\cite[Proposition~1.83]{koch1997algebraic}}]
\label{theorem:kummer_ramification}
Assume $K$ contains the $n$-th roots of unity $\mu_n$ and $(n,p)=1$.
Let $L/K$ be the extension given by $X^n=a$, for $a\in K^\times$.
Then:
\begin{enumerate}
    \item If $v_K(a)\in n\Z$ and the image of
    $a\pi^{-v_K(a)}$ in $\kappa$ is an $n$-th power, the extension is
    trivial.
    \item If $v_K(a)\in n\Z$ and the image of
    $a\pi^{-v_K(a)}$ in $\kappa$ is not an $n$-th power, the extension is
    unramified.
    \item If $v_K(a)\notin n\Z$, the ramification index is
    $n/\gcd(n,v_K(a))$; in particular, if $\gcd(n,v_K(a))=1$, the extension
    is totally and tamely ramified of degree $n$.
\end{enumerate}
\end{theorem}

\subsection{Artin--Schreier theory: order \texorpdfstring{$p$}{p}}

\paragraph{Classes and equations.}
Assume that $K$ has characteristic $p$, and put $\wp(z)=z^p-z$. Write
\[
H^1(K,\Z/p\Z):=H^1_{\mathrm{\acute et}}(\Spec K,\Z/p\Z).
\]
Its elements classify $\Z/p\Z$-torsors over $K$. The Artin--Schreier
sequence gives a natural isomorphism
\[
H^1(K,\Z/p\Z)\cong K/\wp(K).
\]
Under this isomorphism, the coset $[a]\in K/\wp(K)$ corresponds to the
$\Z/p\Z$-torsor
\[
\Spec K[X]/(X^p-X-a),
\]
on which $c\in\F_p\cong\Z/p\Z$ acts by $X\mapsto X+c$. In particular,
replacing $a$ by $a+(b^p-b)$, with $b\in K$, does not change the
isomorphism class of the torsor. If $[a]\neq0$, adjoining a solution of
$X^p-X=a$ gives a cyclic field extension of degree $p$; conversely, every
cyclic extension of degree $p$ arises in this way. When $K=k((x))$ with
$k$ algebraically closed, every class has a representative in
$x^{-1}k[x^{-1}]$ whose nonzero pole order is prime to $p$.

\begin{theorem}[Ramification in Artin--Schreier Extensions,
{\cite[Proposition~2.1]{Thomas2005}}]
\label{theorem:artin_schreier_ramification}
Let $L/K$ be given by $X^p-X=a$. Then:
\begin{enumerate}
    \item If $v_K(a)>0$, or if $v_K(a)=0$ and $a\in\wp(K)$, the extension
    is trivial.
    \item If $v_K(a)=0$ and $a\notin\wp(K)$, the extension is cyclic of
    degree $p$ and unramified.
    \item If $v_K(a)=-m<0$ with $p\nmid m$, the extension is cyclic,
    totally ramified of degree $p$, and has unique upper ramification break
    $m$.
\end{enumerate}
\end{theorem}

\subsection{Artin--Schreier--Witt theory: order
\texorpdfstring{$p^r$}{p-power}}
\phantomsection\label{subsection:witt_vector_prerequisites}

\paragraph{Classes and equations.}
For an $\F_p$-algebra $A$, write $W_r(A)=A^r$ for the ring of $p$-typical
Witt vectors of length $r$, with Witt addition and subtraction denoted by
$\oplus$ and $\ominus$, and put $\wp=F-\mathrm{id}$. For a field $K$ of
characteristic $p$, write
\[
H^1(K,\Z/p^r\Z):=
H^1_{\mathrm{\acute et}}(\Spec K,\Z/p^r\Z).
\]
Its elements classify $\Z/p^r\Z$-torsors over $K$.
Artin--Schreier--Witt theory gives a natural isomorphism
\[
H^1(K,\Z/p^r\Z)\cong W_r(K)/\wp W_r(K).
\]
Under this isomorphism, the coset
$[\vec g]\in W_r(K)/\wp W_r(K)$ corresponds to an affine solution scheme
as follows. Let $\vec X=(X_0,\dots,X_{r-1})$ be the universal Witt vector
on $\bW_{r,K}=\Spec K[X_0,\dots,X_{r-1}]$. Then
\[
T_{\vec g}:=
\Spec\!\left(
K[X_0,\dots,X_{r-1}]
\big/\bigl(\wp\vec X\ominus\vec g\bigr)
\right),
\]
where the parenthesized Witt vector denotes the ideal generated by its $r$
coordinate polynomials. Equivalently, $T_{\vec g}$ is the fiber of
$\wp:\bW_{r,K}\to\bW_{r,K}$ over the $K$-point $\vec g$. Its first defining
equation is $X_0^p-X_0-g_0=0$, and each subsequent equation has the
triangular form
\[
X_j^p-X_j-g_j+
C_j(X_0,\dots,X_{j-1};g_0,\dots,g_{j-1})=0.
\]
Here $C_j$ is the universal Witt correction polynomial determined by the
addition law and depends only on the preceding coordinates. This is a finite
\'etale $K$-scheme of degree $p^r$, on which
$\vec c\in W_r(\F_p)\cong\Z/p^r\Z$ acts by Witt translation
$\vec X\mapsto\vec X\oplus\vec c$. In particular, replacing $\vec g$ by
$\vec g\oplus\wp\vec h$, with $\vec h\in W_r(K)$, gives an isomorphic
torsor. If $[\vec g]$ has order $p^s$, its associated character cuts out a
cyclic field extension of degree $p^s$; thus $T_{\vec g}$ is connected
precisely when $[\vec g]$ has order $p^r$. Conversely, every cyclic
extension of degree dividing $p^r$ arises in this way.
The Witt-vector group scheme is constructed in
\Cref{appendix:witt_vector_algebra_basics}. The Artin--Schreier--Witt
operator, the proof that the displayed equation defines a
$\Z/p^r\Z$-torsor, and the classification theorem---including its description
of the degree and connected components and of induction from a torsor under the
subgroup $\Z/p^s\Z\subseteq\Z/p^r\Z$---are given in
\Cref{appendix:asw_development}.

\begin{definition}\label{definition:reduced_witt_vector}
Let $K=k((x))$. For $g\in K$, put
$m(g):=\max(0,-v_x(g))$. A vector
$\vec g=(g_0,\dots,g_{r-1})\in W_r(K)$ is \textbf{reduced} if, for every
$j$, either $m(g_j)=0$ or $p\nmid m(g_j)$.
(Every class in $W_r(K)/\wp W_r(K)$ admits a reduced representative.)
\end{definition}

\begin{theorem}[Ramification in Artin--Schreier--Witt Extensions
(Schmid--Witt)]
\label{theorem:schmid_witt_break_formula}
Let $K=k((x))$, with $k$ perfect of characteristic $p$, and let $L/K$ be a
cyclic totally ramified extension of degree $p^r$ whose
Artin--Schreier--Witt class is represented by a reduced vector
$\vec g=(g_0,\dots,g_{r-1})$, with $m_j:=m(g_j)$. Then the largest
upper-numbering ramification break of $L/K$ is
\[
u=\max_{0\leq j\leq r-1}p^{r-1-j}m_j.
\]
\end{theorem}

The Schmid--Witt formula is an external ramification theorem: we use it as
stated and do not reproduce its proof. For finite residue fields, the formula
is due to Kanesaka--Sekiguchi \cite{KanesakaSekiguchi1979} and Thomas
\cite{Thomas2005}, building on Schmid \cite{Schmid1937} for $r=1$; for
arbitrary perfect residue fields, it is due to Elder--Keating
\cite{ElderKeating2025}. We use it in one direction only: evaluating the
formula on one reduced representative computes the largest upper break, so
no minimization over representatives is needed.

Reducedness prevents cancellation between Witt levels. Indeed, if two
nonzero pole orders $m_j$ and $m_{j'}$, with $j<j'$, contributed the same
weighted value, then
$m_{j'}=p^{j'-j}m_j$ would be divisible by $p$, contradicting reducedness.
For $r=1$, the formula is exactly the Artin--Schreier break computation.

\subsection{The common realization conclusion}

The three theories give the following single conclusion, which is the form
needed after the formal--local reduction to $\bP^1_k$.

\begin{lemma}[Parallel realization over $\bG_m$]
\label{lemma:realization_over_Gm}
\label{lemma:realization_over_Gm_Wr}
Let $k$ be algebraically closed of characteristic $p>0$, and let $\cQ$ be a
$G$-torsor over $\Spec k((x))$.
\begin{enumerate}
    \item If $\cQ$ is tamely ramified, then there are an integer $e$ prime
    to $p$, an injection $\iota:\Z/e\Z\hookrightarrow G$, and an integer
    $a$ prime to $e$, such that $\cQ$ is the restriction of the induced
    $G$-torsor on $\bG_m$ obtained from
    \[
    T^e=x^a.
    \]
    This torsor is tamely ramified at $0$ and $\infty$.

    \item If $G=\Z/p\Z$ and the ramification of $\cQ$ is bounded by $d$,
    then $\cQ$ is represented by
    \[
    X^p-X=f(x),\qquad f\in x^{-1}k[x^{-1}],
    \]
    where either $f=0$, or
    $m:=\deg_{x^{-1}}f<d$ and $p\nmid m$. The same equation defines a
    torsor on $\bG_m$ extending \'etale across $\infty$.

    \item If $G=\Z/p^r\Z$, the torsor $\cQ$ is connected, and its
    ramification is bounded by $d$, then $\cQ$ is represented by
    \[
    \wp\vec X=\vec f,
    \]
    where $\vec f=(f_0,\dots,f_{r-1})$ is reduced and
    $f_j\in x^{-1}k[x^{-1}]$. Put
    $m_j:=\deg_{x^{-1}}f_j$ when $f_j\neq0$, and $m_j:=0$ when $f_j=0$.
    Then it has the following properties:
    \begin{enumerate}
        \item \label{item:realization_Wr_polebound}
        \begin{equation}\label{eq:parallel_witt_pole_bound}
        p^{r-1-j}m_j<d\qquad(0\leq j<r).
        \end{equation}
        \item \label{item:realization_Wr_torsor}
        The same equation defines a $\Z/p^r\Z$-torsor on $\bG_m$ whose
        restriction to $\Spec k((x))$ is $\cQ$.
        \item \label{item:realization_Wr_infty}
        This torsor extends \'etale across $\infty$.
    \end{enumerate}
\end{enumerate}
\end{lemma}

\begin{proof}
For (1), let
$\rho:\Gal(k((x))^{\mathrm{sep}}/k((x)))\to G$ classify $\cQ$. Since
$k$ is algebraically closed, $k((x))$ has no nontrivial unramified extension
and its tame quotient is procyclic. If $\cQ$ is tame, the image of $\rho$ is
therefore a cyclic group of order $e$ prime to $p$. Let
$\iota:\Z/e\Z\hookrightarrow G$ identify this image with its subgroup of
$G$. Kummer theory identifies the corresponding surjective character with
the class of $x^a$ for some $a$ prime to $e$. Hence $\cQ$ is induced from
$T^e=x^a$. The same equation over $k[x,x^{-1}]$ is \'etale and defines the
required torsor on $\bG_m$; its ramification at both boundary points is tame.

For (2), Artin--Schreier theory identifies $\cQ$ with a class in
$k((x))/\wp(k((x)))$. The length-one construction in the proof of
\Cref{prop:reduced_polynomial_representatives} produces a representative
$f\in x^{-1}k[x^{-1}]$ which is zero or has pole order $m$ prime to $p$.
By \Cref{theorem:artin_schreier_ramification}, the ramification bound is
equivalent to $m<d$. The equation $X^p-X=f$ is \'etale over
$k[x,x^{-1}]$; since $f\in k[x^{-1}]$, it also extends over
$\Spec k[x^{-1}]=\bP^1_k\setminus\{0\}$.

For (3), let $[\vec f_0]\in W_r(k((x)))/\wp$ be the class of $\cQ$ and
choose a reduced polynomial representative $\vec f$ using
\Cref{prop:reduced_polynomial_representatives}. Connectedness and
\Cref{cor:asw_class_order_and_first_component} imply $f_0\neq0$. We use
the convention for $m_j$ in the statement. By
\Cref{theorem:schmid_witt_break_formula}, the Schmid--Witt formula gives
\[
\max_j p^{r-1-j}m_j<d,
\]
which is \Cref{eq:parallel_witt_pole_bound}. By
\Cref{lemma:asw_etale_torsor}, the equation
$\wp\vec X=\vec f$ defines a $\Z/p^r\Z$-torsor over $\bG_m$ with the
prescribed local class. Since every $f_j$ lies in $k[x^{-1}]$, the same lemma
applied over $k[x^{-1}]$ extends the torsor \'etale across $\infty$.
\end{proof}

\begin{rem*}\label{rem:realization_Wr_sanity}
In part~(3), the weighted pole bound forces $f_j=0$ whenever
$p^{r-1-j}\geq d$: if $f_j\neq0$, then $m_j$ is a positive integer and
$p^{r-1-j}m_j\geq d$, contradicting
\Cref{eq:parallel_witt_pole_bound}. Thus a break smaller than $d$ imposes a
visible restriction on every Witt component, not only on the last one.
\end{rem*}

\subsection{Three common inputs for the cyclic calculations}

The explicit equations above enter the symmetric-power construction through
one contracted-product calculation. We record it once for all three cyclic
theories.

Let $S$ be a scheme, let $H$ be a commutative $S$-group scheme, and let
$\varphi:H\to H$ be a finite \'etale surjective homomorphism with constant
kernel $G$. For a section $a:X\to H$, write
\[
    \mathcal T_a:=X\times_{a,H,\varphi}H.
\]
Thus $\mathcal T_a$ is a $G$-torsor on $X$. We write $\star$ for the group
law of $H$: multiplication on $\bG_m$, ordinary addition on $\bG_a$, and
Witt addition on $\bW_r$.

\begin{prop}[Common cyclic product calculation]
\label{prop:parallel_cyclic_product_torsor}
Let $a_i:X\to H$, $1\leq i\leq d$, be sections. There is a canonical
isomorphism of $G$-torsors
\[
    \mathcal T_{a_1}\otimes^G\cdots\otimes^G\mathcal T_{a_d}
       \cong \mathcal T_{a_1\star\cdots\star a_d}.
\]
On equation coordinates it is induced by
\[
    (X_1,\dots,X_d)\longmapsto
    Y:=X_1\star\cdots\star X_d,
    \qquad
    \varphi(Y)=a_1\star\cdots\star a_d.
\]
If a group $\Gamma$ permutes the indices, this isomorphism is
$\Gamma$-equivariant and therefore compatible with descent to the quotient.
\end{prop}

\begin{proof}
The product $\prod_i\mathcal T_{a_i}$ is a $G^d$-torsor. Contracting along
the multiplication homomorphism $G^d\to G$ gives its quotient by $\KG$.
The displayed map is $\KG$-invariant because the kernel translations cancel
under $\star$, and it is equivariant for the residual $G$-action. Its target
is a $G$-torsor because it is the pullback of $\varphi$. Hence it is an
isomorphism by \Cref{prop:torsor_morphism_iso}. Commutativity and
associativity of $\star$ give the permutation-equivariance.
\end{proof}

We shall also pass from the connected inertia torsor to the original torsor.
The following compatibility ensures that this passage commutes with the
symmetric-power construction.

\begin{lemma}[Symmetric powers and induced torsors]
\label{lemma:symmetric_power_induced_torsor}
Let $\iota:H\hookrightarrow G$ be an injection of finite abelian groups, and
let $\cQ$ be an $H$-torsor on $U$. There is a canonical isomorphism of
$G$-torsors over $U^{(d)}$
\[
    (\iota_*\cQ)^{[d]}\cong\iota_*(\cQ^{[d]}).
\]
\end{lemma}

\begin{proof}
The canonical morphism
$j:\cQ\to\iota_*\cQ=(\cQ\times G)/H$ is $H$-equivariant, where $H$ acts
on the target through $\iota$. Its $d$-fold product commutes with the
$S_d$-actions and is equivariant for $H^d\to G^d$. Writing
$K_H:=\ker(H^d\to H)$, it therefore descends to an $H$-equivariant morphism
\[
    \cQ^{[d]}=(K_H\rtimes S_d)\backslash\cQ^{\times d}
       \longrightarrow
    (\KG\rtimes S_d)\backslash(\iota_*\cQ)^{\times d}
       =(\iota_*\cQ)^{[d]}.
\]
The universal property of induction extends this morphism uniquely to a
$G$-equivariant morphism
$\iota_*(\cQ^{[d]})\to(\iota_*\cQ)^{[d]}$. Both sides are $G$-torsors over
$U^{(d)}$, so it is an isomorphism by
\Cref{prop:torsor_morphism_iso}.
\end{proof}

Finally, the two wild calculations use the same elementary fact about the
exceptional valuation. Let $e_1,\dots,e_d$ be the elementary symmetric
polynomials in $x_1,\dots,x_d$, put $y_i=x_i^{-1}$, and write
\[
    \kappa:=k(e_1/e_d,\dots,e_{d-1}/e_d).
\]

\begin{lemma}[Symmetric regularity]
\label{lemma:symmetric_regularity}
If $\Phi\in k[y_1,\dots,y_d]^{S_d}$ and every monomial of $\Phi$ has total
degree strictly less than $d$, then
\[
    \Phi\in k(e_1/e_d,\dots,e_{d-1}/e_d)=\kappa.
\]
In particular, $\Phi$ is regular at the generic point of the exceptional
divisor.
\end{lemma}

\begin{proof}
The identity
\[
    e_j(y_1,\dots,y_d)
       =\frac{e_{d-j}(x_1,\dots,x_d)}{e_d(x_1,\dots,x_d)}
\]
puts $e_1(y),\dots,e_{d-1}(y)$ in $\kappa$. By the fundamental theorem of
symmetric polynomials, each homogeneous part of $\Phi$ is a polynomial in
$e_1(y),\dots,e_d(y)$ with their natural degrees. Since every homogeneous
part has degree $<d$, the generator $e_d(y)$ cannot occur. Hence
$\Phi\in\kappa\subset\kappa[[e_d]]$.
\end{proof}
